Nell'ultimo decennio, la crescente automazione dei processi produttivi ha trasformato radicalmente le aspettative nei confronti dei sistemi di controllo qualità. Individuare un difetto su un componente meccanico, riconoscere un'imperfezione su una superficie verniciata, rilevare un'anomalia in un circuito stampato: operazioni che un tempo richiedevano un operatore umano vengono oggi delegate, con risultati sempre più affidabili, a sistemi di visione artificiale basati su algoritmi di deep learning. Al centro di questa trasformazione si colloca il problema del \textit{visual anomaly detection}, ovvero il rilevamento automatico di pattern visivi inattesi in immagini di prodotti industriali.

L'anomaly detection industriale si distingue dalla classificazione convenzionale per una caratteristica fondamentale: durante l'addestramento, il sistema dispone esclusivamente di esempi normali, poiché i difetti sono rari, costosi da raccogliere e, per natura, imprevedibili nella loro manifestazione. Questa situazione ha portato allo sviluppo di una famiglia eterogenea di approcci, ciascuno con assunzioni, punti di forza e limitazioni proprie: dai metodi basati su autoencoder e reti generative avversarie, fino ai più recenti modelli student-teacher, alle architetture memory-based e ai vision-language model. 

In parallelo, sono stati definiti dataset standardizzati che si sono imposti come benchmark, rendendo possibile misurare con precisione i progressi del campo e indirizzare lo sviluppo verso le sfide ancora aperte.

L'obiettivo di questa tesi è offrire un'analisi comparativa dell'evoluzione dei metodi e dei dataset per il visual anomaly detection in ambito industriale, tracciando un percorso che parte dalle fondamenta teoriche dell'apprendimento automatico e arriva agli sviluppi più recenti del 2025. Il lavoro si propone non solo di catalogare i contributi della letteratura, ma di metterli in relazione tra loro, evidenziando le linee evolutive, le discontinuità paradigmatiche e le tendenze ancora in corso.

La trattazione è organizzata in tre capitoli.

\medskip
\noindent Il \textbf{Capitolo 1} introduce i fondamenti teorici necessari alla comprensione del resto della tesi: si parte dalla definizione formale di machine learning e dei suoi paradigmi principali (apprendimento supervisionato, non supervisionato e per rinforzo) per poi approfondire i concetti di representation learning, reti neurali feedforward e ricorrenti, algoritmi di training e tecniche di regolarizzazione. Il capitolo si conclude con un'analisi delle reti neurali convoluzionali (CNN), esaminando i principi fondamentali (connettività locale, condivisione dei pesi e pooling), l'evoluzione storica delle principali architetture e il ruolo del transfer learning, tecnica che si rivelerà centrale in quasi tutti i metodi di anomaly detection discussi nei capitoli successivi.

\medskip
\noindent Il \textbf{Capitolo 2} costituisce il nucleo della tesi. Dopo aver formalizzato il problema dell'anomaly detection e discusso le sue specificità nel contesto industriale, viene ripercorsa cronologicamente l'evoluzione dei metodi di deep learning per il rilevamento di anomalie visive. L'analisi prende le mosse dai primi approcci basati su reti generative avversarie (AnoGAN, GANomaly, f-AnoGAN) e dagli autoencoder deterministici e variazionali, per poi esaminare la svolta rappresentata dai metodi feature-based (PaDiM, SPADE, Uninformed Students), i modelli a normalizing flow (CFLOW-AD, FastFlow, CS-Flow), le tecniche di generazione sintetica di anomalie (CutPaste, DRAEM), i metodi memory-based (PatchCore, MemSeg), le architetture student-teacher avanzate (Reverse Distillation, RD++, EfficientAD) e i metodi semplificati come SimpleNet. 

Il capitolo dedica infine una sezione ai vision-language model (WinCLIP, AnomalyCLIP) e all'anomaly detection su dati tridimensionali, con i metodi feature-based per point cloud (BTF, Shape-Guided, Reg3D-AD) e i framework multimodali RGB+3D (AST, M3DM, CFM, D\textsuperscript{3}M). 

Ciascuna famiglia di metodi è analizzata nelle sue motivazioni, architettura, vantaggi e limitazioni, con particolare attenzione alle discontinuità che hanno definito il passaggio da un paradigma al successivo.

\medskip
\noindent Il \textbf{Capitolo 3} sposta il focus dai metodi ai dati, esaminando l'evoluzione dei principali benchmark per l'anomaly detection industriale dal 2018 al 2025. Vengono analizzati in dettaglio quattordici dataset, mettendone in evidenza le scelte di design, le categorie di oggetti, le tipologie di difetti e le metriche di valutazione adottate. 

L'analisi mostra come l'orizzonte delle sfide proposte si sia ampliato in maniera crescente: dalla detection di difetti superficiali su texture strutturate, fino a scenari che richiedono ragionamento logico e strutturale \cite{bergmann2022beyond}, acquisizioni 3D \cite{bergmann2021mvtec3d}, modalità multiple \cite{zhu2025realiadd3, sim3d2025} e condizioni più vicine alla produzione reale \cite{wang2024realiad, bergmann2025mvtecad2}.

\medskip
Nel complesso, il panorama che emerge è quello di un campo in rapida crescita, in cui metodi sempre più efficienti e generalizzabili si confrontano con benchmark sempre più ricchi e realistici. L'auspicio è che questa analisi possa costituire un punto di riferimento utile per chiunque si avvicini al visual anomaly detection industriale, facilitando l'orientamento in una letteratura densa e in continua evoluzione.
